\Needspace{18\baselineskip}
\section{State and contextual consequences}

The following results compose the two simulation constructions with the
soundness and observation lemmas. They are short in Lean because their
mathematical work has been isolated in those earlier results.

\begin{theorem}[Converse state sufficiency]
If $M$ is deterministic, $(Q,P,S)$ is reasonable,
$\safe{Q^{-1}}{P}$ holds, and $G^{\mathrm{op}}$ is robust for
$(Q,P,S)$, then
\[
 \Lft(s,t)\ \Longrightarrow\ t\preceq_Q s.
\]
\end{theorem}
\lean{LocalReasoning}{281}{286}{SpecimenSystem.sufficiency_converse}
\begin{proof}
The lifting-opposite lemma changes the premise into
$\mathcal L_{G^{\mathrm{op}}}(t,s)$. The converse construction supplies
a counting simulation on this relation. Apply counting-simulation
soundness to obtain the stated refinement.
\end{proof}

\begin{theorem}[Mutual state refinement]
Suppose $M$ is deterministic, $(Q,P,S)$ is reasonable, $G$ is
input-safe for both $(Q,P)$ and $(Q^{-1},P)$, and both $G$ and
$G^{\mathrm{op}}$ are robust for $(Q,P,S)$. Then
\[
 \Lft(s,t)\ \Longrightarrow\ s\simeq_Q t.
\]
\end{theorem}
\lean{LocalReasoning}{289}{295}{SpecimenSystem.equivalence}
\begin{proof}
Forward state sufficiency supplies $s\preceq_Q t$. Converse state
sufficiency supplies $t\preceq_Q s$. These are exactly the two
components of state equivalence.
\end{proof}
The main cost relation $Q$ itself need not be symmetric. The two
directional refinement statements use the same $Q$, with source and
target counts exchanged when the states are exchanged.

For the contextual results, fix the observation interface of the earlier
section and a family $T_i$ of relations on fragments. Assume the
following \emph{initial-lifting property}: for $T_i(a,b)$, every
allowed context $c$ and observations
$\mathsf{Obs}(c,a,s)$, $\mathsf{Obs}(c,b,t)$ give $\Lft(s,t)$.
This property is the explicit premise \texttt{initial\_lift} in the
source. It connects the abstract specimen system with the separately
supplied observation relation.

\begin{theorem}[Forward contextual sufficiency]
If $M$ is deterministic, $(Q,P,S)$ is reasonable,
$\safe{Q}{P}$ and $\robust{Q}{P}{S}$ hold, and the initial-lifting
property holds, then
\[
 T_i(a,b)\ \Longrightarrow\ a\preceq_Q^{\mathcal O}b.
\]
\end{theorem}
\lean{LocalReasoning}{311}{316}{SpecimenSystem.contextual_sufficiency}
\begin{proof}
Use contextual refinement from a lifted relation, taking $L=\Lft$.
Its initial-lifting premise is precisely the one assumed here. Its
state-soundness premise is forward state sufficiency. Applying that
proposition to $T_i(a,b)$ gives the desired contextual refinement.
\end{proof}

\begin{theorem}[Converse contextual sufficiency]
If $M$ is deterministic, $(Q,P,S)$ is reasonable,
$\safe{Q^{-1}}{P}$ holds, $G^{\mathrm{op}}$ is robust for
$(Q,P,S)$, and the same initial-lifting property holds, then
\[
 T_i(a,b)\ \Longrightarrow\ b\preceq_Q^{\mathcal O}a.
\]
\end{theorem}
\lean{LocalReasoning}{318}{323}{SpecimenSystem.contextual_sufficiency_converse}
\begin{proof}
Fix $T_i(a,b)$, an allowed context, an observation of $b$ producing
$s$, and an observation of $a$ producing $t$. Supply the observations
to initial lifting in the original $a,b$ order to get $\Lft(t,s)$.
Converse state sufficiency then gives $s\preceq_Q t$, the required
contextual comparison.
\end{proof}

These are conditional abstract counterparts of the two clauses of the
paper's Theorem~6. Refocusing, output-closure, type-correct plugging, and
the paper's concrete specimen definitions have not disappeared from the
mathematics: their required consequences are currently represented by
the specimen-system laws and initial-lifting premise. Instantiating
those hypotheses is the remaining concrete formalization problem.
